\chapter{Syntax highlighting}
\label{ch:syntax}

\begin{aside}
A code viewer without syntax highlighting feels primitive.
Coloured keywords, strings, comments --- the cost is a
parser; the benefit is readability.
\end{aside}

\section{Approaches}

\begin{itemize}[leftmargin=*]
  \item \textbf{Regex.} Simple, fast, imperfect.
  \item \textbf{Tree-sitter.} Real parser, incremental, accurate.
  \item \textbf{LSP semantic tokens.} From a language server,
    most accurate.
\end{itemize}

\maya{} doesn't bundle a highlighter. Pick by need.

\section{Tree-sitter}

Tree-sitter parses on every edit incrementally; only the
changed subtree re-parses. Cost is O(edit), not O(file).

\begin{lstlisting}
auto tree = ts_parser_parse_string(parser, nullptr, source, len);
auto root = ts_tree_root_node(tree);
// walk nodes, map to styles
\end{lstlisting}

\section{Mapping to styles}

A theme assigns styles to syntax categories:

\begin{lstlisting}
[syntax]
keyword = { fg = "purple", attrs = ["bold"] }
string = { fg = "green" }
comment = { fg = "grey", attrs = ["italic"] }
\end{lstlisting}

\section{Incremental rendering}

Highlight only the visible viewport. As the user scrolls,
highlight new rows on demand.

\section{Cache invalidation}

On edit, invalidate the highlight cache for the changed
region. Use tree-sitter's incremental parsing to limit
the work.

\section{Languages}

A real editor supports dozens. Each language ships a
grammar; load on demand.

\section{Recap}

\begin{recap}
\begin{itemize}[leftmargin=*]
  \item Regex for simple cases; tree-sitter for real.
  \item Theme maps syntax categories to styles.
  \item Highlight only the visible region.
  \item Incremental on edit.
\end{itemize}
\end{recap}

\section*{Workshop}

\begin{enumerate}[leftmargin=*]
  \item Add basic regex highlighting (keyword, string,
    comment) to your code viewer.
  \item Integrate tree-sitter for one language.
  \item Build a theme switcher for syntax styles.
\end{enumerate}
